\chapter{1996:Robert F. Curl Jr. and Sir Harold Kroto and Richard E. Smalley}

\label{chap:chemistry1996}

The Nobel Prize in Chemistry for the year 1996 was awarded jointly to Robert F. Curl Jr. and Sir Harold Kroto and Richard E. Smalley.

\textbf{Motivation:}

for their discovery of fullerenes

\section{Robert F. Curl Jr.}

\subsection*{Early Life and Education}

Robert F. Curl Jr. was born in 1933 in Alice, Texas, USA.

\subsection*{Career and Major Research}

Curl studied at Rice University (B.A., 1954) and earned his Ph.D. from the University of California, Berkeley, in 1957. He joined the faculty of Rice University in 1958, where he spent his career.

\subsection*{Nobel Prize-winning Work}

The work for which Robert F. Curl Jr. was awarded the Nobel Prize in Chemistry in 1996 was described by the Nobel Committee as: "for their discovery of fullerenes".

In experiments with Kroto and Smalley, laser vaporization of graphite produced clusters of carbon, among which the 60-atom molecule C60 was identified as a truncated icosahedron.

\subsection*{Significance and Impact}

The discovery of buckminsterfullerene opened the field of fullerene chemistry and nanostructured carbon materials.

\subsection*{Later Career and Honors}

Curl was a professor emeritus at Rice University. He shared the 1996 Nobel Prize in Chemistry, and he died in 2022.

\subsection*{Legacy}

Fullerenes became a new branch of carbon chemistry and materials science.

\subsection*{Selected Publications}

"C60: Buckminsterfullerene" (Nature, 1985)

\clearpage

\section{Sir Harold Kroto}

\subsection*{Early Life and Education}

Sir Harold Kroto was born in 1939 in Wisbech, England.

\subsection*{Career and Major Research}

Kroto studied at the University of Sheffield (B.Sc., 1961; Ph.D., 1964). He conducted research at the University of Sussex for most of his career and later held a position at Florida State University.

\subsection*{Nobel Prize-winning Work}

The work for which Sir Harold Kroto was awarded the Nobel Prize in Chemistry in 1996 was described by the Nobel Committee as: "for their discovery of fullerenes".

With Curl and Smalley he identified C60, whose hollow cage structure was suggested by the geodesic domes designed by Buckminster Fuller.

\subsection*{Significance and Impact}

The identification of C60 revealed an entirely new molecular form of carbon.

\subsection*{Later Career and Honors}

Kroto was knighted for services to science and education. He shared the 1996 Nobel Prize in Chemistry, and he died in 2016.

\subsection*{Legacy}

His advocacy for science education, and the discovery itself, left a lasting mark on chemistry and science communication.

\subsection*{Selected Publications}

"C60: Buckminsterfullerene" (Nature, 1985)

\clearpage

\section{Richard E. Smalley}

\subsection*{Early Life and Education}

Richard E. Smalley was born in 1943 in Akron, Ohio, USA.

\subsection*{Career and Major Research}

Smalley studied at the University of Michigan (B.S., 1965) and earned his Ph.D. from Princeton University in 1973. He joined Rice University in 1976.

\subsection*{Nobel Prize-winning Work}

The work for which Richard E. Smalley was awarded the Nobel Prize in Chemistry in 1996 was described by the Nobel Committee as: "for their discovery of fullerenes".

He built the laser-vaporization apparatus in which fullerenes were discovered and helped establish the structure of C60.

\subsection*{Significance and Impact}

The discovery demonstrated that carbon forms stable hollow cage molecules, launching the field of nanoscale carbon science.

\subsection*{Later Career and Honors}

Smalley later directed the Center for Nanoscale Science and Technology at Rice University. He shared the 1996 Nobel Prize in Chemistry, and he died in 2005.

\subsection*{Legacy}

His vision of carbon nanotubes and nanoscale materials helped establish nanotechnology as a research field.

\subsection*{Selected Publications}

"C60: Buckminsterfullerene" (Nature, 1985)

\clearpage

\section*{Chapter Summary}

The 1996 Nobel Prize in Chemistry recognized the discovery of fullerenes. Curl, Kroto, and Smalley found that carbon can form hollow cage molecules such as C60, opening up a new field of carbon chemistry.
